%% 共通プリアンブル。main.tex（本体）と solutions.tex（別冊解答）の両方が読む。
%% corona-a5-1.1.cls は改変禁止なので，上書きはすべてここに置く。

\usepackage[dvips]{graphicx}
\usepackage[fleqn]{amsmath}
\usepackage{amssymb}
\usepackage{type1cm}
\usepackage{multicol}
\usepackage{makeidx}

\usepackage{alltt}
\usepackage{ascmac}

\usepackage{a5update}
\usepackage{a5update-fix01}
\usepackage{linequote}

\usepackage{fix_for_tl2021}

\usepackage{lawenv}     %% 「法則」環境（公理＝出発点の基本法則用）
\usepackage{etoolbox}   %% \patchcmd 用（下記の \@startsection 修正）

%%====================================================================
%% 組版の修正（2026-09-02 Fableレビュー C群対応）
%% corona-a5-1.1.cls は改変禁止のため，必要な上書きはここで行う。
%%====================================================================
\makeatletter

%% [C-1] 問環境の直後の \section で必ず改頁されるのを直す
%%   a5update.sty の 問 環境は末尾で \nointerlineskip\addvspace{20H} を行う
%%   （\prevdepth=-1000pt のまま終わる）。一方 cls の \@startsection は
%%   レベル1見出しで \vskip-\prevdepth を入れるため 1000pt の空白が生じ，
%%   強制改頁になっていた。\prevdepth が -1000pt のときは代わりに直前の
%%   後アキ（\addvspace 分）を打ち消し，通常の見出し前アキだけを残す。
\patchcmd{\@startsection}%
  {\vskip-\prevdepth \prevdepth\z@}%
  {\ifdim\prevdepth>-\@m\p@ \vskip-\prevdepth \else \vskip-\lastskip \fi
   \prevdepth\z@}%
  {}{\ClassError{main}{\string\@startsection\space の修正に失敗}{}}

%% [C-4] COLUMN 環境に題名のオプション引数を追加
%%   cls の COLUMN は引数を取らず，題名は本文先頭に \FIVEjidori{{\bfseries 題名}}
%%   を書く様式（電磁気本と同じ）。\begin{COLUMN}[題名] でも同じ出力になるよう
%%   上書きする。オプションを省略すれば従来どおり。
\let\pe@COLUMN\COLUMN
\let\pe@endCOLUMN\endCOLUMN
\renewenvironment{COLUMN}[1][]{%
  \pe@COLUMN
  \if\relax\detokenize{#1}\relax\else\FIVEjidori{{\bfseries #1}}\par\fi}%
  {\pe@endCOLUMN}

%% [C-5] 序章（\chapter*，章番号0）の図表番号を「図0.1」「表0.1」にする
%%   cls の \thefigure は章番号0のとき「図1」と出すため，序章の前後で差し替える。
\newcommand{\PEprologuenumbering}{%
  \let\pe@thefigure\thefigure
  \let\pe@thetable\thetable
  \renewcommand{\thefigure}{図0.\@arabic\c@figure}%
  \renewcommand{\thetable}{表0.\@arabic\c@table}%
  \setcounter{figure}{0}\setcounter{table}{0}}
\newcommand{\PErestorenumbering}{%
  \let\thefigure\pe@thefigure
  \let\thetable\pe@thetable}

\makeatother

%% 校正用：修正箇所を赤字で示す（2026-09-02 Fableレビュー対応）
%% 最終版では \revcolorfalse に切り替えれば全て黒に戻る（原稿はそのまま）。
%% color は corona-a5-1.1.cls が [usenames] 付きで読込済み（重複読込するとoption clash）
\newif\ifrevcolor
\revcolortrue
\newcommand{\rev}[1]{\ifrevcolor\textcolor{red}{#1}\else#1\fi}

%% a5update.sty が amsmath の \tagform@ を再定義し \eqref が壊れるため上書き
%% （この環境では \ref が式番号を括弧付きで返す）
\renewcommand*{\eqref}[1]{\ref{#1}}

